% ==============================================================================
% Plantilla Institucional de Trabajo Terminal — UPIIZ IPN
% Archivo: capitulos/09-conclusiones.tex
% ==============================================================================

\chapter{Resultados Esperados y Conclusiones del Protocolo}
\label{cap:conclusiones}

Con base en la formulación analítica, el plan de trabajo propuesto y la articulación sinérgica de las cuatro áreas de la ingeniería mecatrónica, se proyecta alcanzar los siguientes resultados y contribuciones al concluir las etapas de Trabajo Terminal I y II:

\begin{enumerate}
    \item \textbf{Proyección respecto al Objetivo Específico 1:} Se obtendrá una representación dinámica continua y discreta en espacio de estados que modele con fidelidad el comportamiento acoplado del tren de tracción y del mecanismo de dirección Ackermann, sustentada en la identificación paramétrica experimental de los componentes electromecánicos.
    \item \textbf{Proyección respecto al Objetivo Específico 2:} Se demostrará la controlabilidad y observabilidad del sistema mecatrónico, sintetizando la matriz de ganancias de realimentación $\mathbf{K}$ que garantizará una respuesta transitoria rápida ($t_s \le \qty{0.50}{\second}$) con un sobrepaso controlado ($M_p \le 10\,\%$) y sin saturación excesiva de las etapas de potencia.
    \item \textbf{Proyección respecto al Objetivo Específico 3:} Se consolidará una plataforma vehicular a escala 1:10 instrumentada, modular y reproducible, con electrónica de potencia aislada y firmware determinista en lenguaje C++ ejecutado a \qty{1}{\kilo\hertz} sobre un microcontrolador de 32 bits.
    \item \textbf{Proyección respecto al Objetivo Específico 4:} Se establecerá un protocolo de validación experimental cuantificable que permitirá verificar en pista el cumplimiento de los márgenes de error especificados ($\text{MAE}_v \le \qty{0.05}{\meter\per\second}$, $\text{RMSE}_v \le \qty{0.08}{\meter\per\second}$), evaluando la capacidad de rechazo ante perturbaciones y la robustez del sistema ante caídas de tensión de la batería.
\end{enumerate}

En conclusión, la ejecución del proyecto bajo los lineamientos de la Línea IV de la UPIIZ-IPN proporcionará un banco de pruebas mecatrónico de alto valor académico, sirviendo de base experimental sólida para futuras investigaciones en sistemas de control y vehículos inteligentes.
